\documentclass[12pt]{article}
\usepackage[T1]{fontenc}
\usepackage[utf8]{inputenc}
\usepackage{lmodern}
\usepackage{textcomp}
\usepackage{enumitem}
\usepackage{geometry}
\geometry{margin=1in}
\begin{document}
\begin{center}
Answer Key - Philosophy - Graduate Level Assessment 9 \
\small (Answers are ordered exactly as in the question paper.)
\end{center}

\section*{Multiple Choice}
\begin{enumerate}[label=\arabic*.]
\item \textbf{Answer:} A
\\ \textit{Options:} A) the greatest harm that one can suffer.; B) an illusion that does not truly exist.; C) the transformation into a higher being.; D) one component of the continuous cycle of rebirth.; E) a state of eternal suffering.
\\ \textit{Note:} Epicurus views death as the greatest harm because he argues that it is the cessation of consciousness and pleasure, leading to the ultimate loss of experiencing life and the good, which he considers the highest goal of human existence.
\item \textbf{Answer:} D
\\ \textit{Options:} A) justice; B) fairness; C) respect for other persons; D) all of the above
\\ \textit{Note:} Arthur argues that values such as well-being, equality, and community are more fundamental than rights and desert, as they provide a deeper ethical foundation for evaluating moral claims and social justice.
\item \textbf{Answer:} B
\\ \textit{Options:} A) a categorical enthymeme; B) a conditional syllogism; C) a disjunctive inference; D) a distributive syllogism; E) a hypothetical syllogism
\\ \textit{Note:} The statement is correct because a conditional syllogism, which is structured around an "if-then" premise, requires that the minor premise either affirms the antecedent of the conditional statement or negates the consequent to maintain logical validity.
\item \textbf{Answer:} A
\\ \textit{Options:} A) with a focus on quantity over quality, leading to possible misallocation.; B) very inefficiently, and are therefore not worth donating to.; C) with a high level of bureaucracy and corruption, therefore not recommended.; D) more effectively than we could get it to our close neighbors.; E) extremely efficiently, making them the only viable option for aid distribution.
\\ \textit{Note:} Singer argues that famine relief organizations prioritize the sheer number of individuals aided rather than the effectiveness or quality of the assistance provided, which can result in resources being misallocated and not addressing the most pressing needs of those affected.
\item \textbf{Answer:} D
\\ \textit{Options:} A) appeal to ignorance; B) straw man; C) bandwagon; D) division; E) false cause
\\ \textit{Note:} The fallacy of division occurs when one assumes that what is true of a whole must also be true of its individual parts, which is a logical error in reasoning.
\end{enumerate}

\section*{True / False}
\begin{enumerate}[label=\arabic*.]
\setcounter{enumi}{5}
\item \textbf{Answer:} True
\\ \textit{Options:} A) True; B) False
\\ \textit{Note:} According to Kant's deontological ethics, the moral worth of an action is determined by the principle or maxim that underlies the action, rather than the consequences of the action itself, as he emphasizes the importance of acting according to duty derived from rational moral laws.
\item \textbf{Answer:} False
\\ \textit{Options:} A) True; B) False
\\ \textit{Note:} Singer argues that the imposition of Western educational systems often undermines rather than enhances the cultural identity of poorer nations by promoting Western values over indigenous practices and knowledge.
\item \textbf{Answer:} False
\\ \textit{Options:} A) True; B) False
\\ \textit{Note:} The principle of military necessity, which justifies the use of force to achieve legitimate military objectives while adhering to ethical constraints, is not inherently aligned with moral nihilism, which denies any objective moral values, but rather emphasizes the importance of moral considerations in the conduct of war.
\item \textbf{Answer:} False
\\ \textit{Options:} A) True; B) False
\\ \textit{Note:} A sarcastic response may express disdain or ridicule but does not necessarily constitute an argumentum ad hominem fallacy unless it explicitly attacks the opponent's character in place of addressing the argument itself.
\item \textbf{Answer:} True
\\ \textit{Options:} A) True; B) False
\\ \textit{Note:} Hume argues that reason is inherently linked to our passions, serving as a tool to inform and guide our emotional responses rather than existing as an independent force, thereby emphasizing the interdependence between reason and passion in human motivation and decision-making.
\end{enumerate}

\section*{Long Form Response}
\begin{enumerate}[label=\arabic*.]
\setcounter{enumi}{10}
\item \textbf{Answer:} To construct the truth table for the argument (G ≡ H) · \textasciitilde{}I, \textasciitilde{}G ∨ (\textasciitilde{}H ∨ I) / G, we analyze the truth values of the premises and the conclusion. The argument is invalid, as demonstrated by a counterexample: when G, H, and I are all false. In this scenario, (G ≡ H) evaluates to true (since both are false), \textasciitilde{}I is true, making the conjunction (G ≡ H) · \textasciitilde{}I true. The disjunction \textasciitilde{}G ∨ (\textasciitilde{}H ∨ I) is also true because \textasciitilde{}G is true. However, the conclusion G is false, thus failing to uphold the validity of the argument. This highlights that even if the premises are satisfied, the conclusion does not necessarily follow, indicating the argument's invalidity.
\\ \textit{Note:} correct because it accurately constructs the truth table for the argument, demonstrates that the premises can be true while the conclusion is false, thereby proving the argument's invalidity through a specific counterexample.
\item \textbf{Answer:} Nathanson argues that the act of murder fundamentally alters the moral landscape of an individual's rights, particularly the right to privacy. He posits that committing such an egregious act constitutes a forfeiture of privacy rights because it symbolizes a rejection of the ethical principles that underpin those rights. This perspective raises significant philosophical implications, suggesting that moral accountability can lead to a reassessment of individual rights in the face of heinous actions. Consequently, the relationship between moral culpability and the forfeiture of rights prompts a deeper inquiry into the balance between societal protection and individual autonomy. Ultimately, Nathanson's argument challenges the notion of privacy as an inviolable right, especially when weighed against the gravity of moral transgressions like murder.
\\ \textit{Note:} correct because it accurately reflects Nathanson's argument that murder undermines the ethical foundation of privacy rights, leading to the conclusion that individuals who commit such acts forfeit their right to privacy, thereby raising important philosophical questions about moral accountability and the reassessment of rights.
\end{enumerate}

\vfill
\noindent\textit{End of Answer Key}
\end{document}